\documentclass[a4paper, 12pt]{article}

% Корректное отображение шрифтов и символов
\usepackage[T2A]{fontenc} 
\usepackage[english, russian]{babel}
\usepackage[utf8]{inputenc}
\usepackage{amsmath}
\usepackage{amssymb}
\usepackage{amsthm}
\usepackage{mathtools}

\usepackage[left=2cm,right=1.5cm,top=2cm,bottom=2cm]{geometry}

% Корректная работа с кириллицей (чтобы работал Ctrl+F)
\usepackage{cmap}

% Картинки и графики
\usepackage{graphicx}
\usepackage{tikz}
% \usepackage{asymptote}

% Удобная команда для отступов
\newcommand\tab[1][.5cm]{\hspace*{#1}}

% Для удобства
\newcommand{\Z}{\mathbb{Z}}
\newcommand{\N}{\mathbb{N}}
\newcommand{\R}{\mathbb{R}}
\newcommand{\Q}{\mathbb{Q}}
\renewcommand{\div}{\mid}
\newcommand{\ndiv}{\nmid}

% Красивые phi, epsilon и пустое множетсво
\renewcommand{\phi}{\varphi}
\renewcommand{\epsilon}{\varepsilon}
\renewcommand{\emptyset}{\varnothing}

% Было использовано по неопытности:
\newcommand{\lra}{\Leftrightarrow}
\newcommand{\om}{\bar{\bar{o}}}

% Математические окружения:
% С счетчиками
\theoremstyle{definition}
\newtheorem{definition}{Определение}[section]
\newtheorem{theorem}{Теорема}[section]
\newtheorem{cons}{Следствие}
\newtheorem{lemma}{Лемма}[section]
\newtheorem{statement}{Утверждение}[section]
% Без счетчиков
\newtheorem*{consequense}{Следствие}
\newtheorem*{comm}{Замечание}
\newtheorem*{axiom}{Аксиома}
\newtheorem*{example}{Пример}
\newtheorem*{examples}{Примеры}
\newtheorem*{theorem*}{Теорема}
% Пришлось добавить
\newtheorem*{introduction}{Введение}
\newtheorem*{algorithm}{Алгоритм}
\newtheorem*{eratosphene}{Решето Эратосфена}
\newtheorem*{def20}{Определение 2.0}

% Настройка заголовков
\usepackage{titlesec}
\titleformat{\section}{\LARGE \bfseries}{\thesection}{1em}{}
\titleformat{\subsection}{\Large\bfseries}{\thesubsection}{1em}{}
\titleformat{\subsubsection}{\large\bfseries}{\thesubsubsection}{1em}{}

% Настройка гиперссылок
\usepackage{hyperref}
\usepackage{xcolor}
\definecolor{linkcolor}{HTML}{225ae2}
\definecolor{urlcolor}{HTML}{225ae2}
\hypersetup{
    pdfstartview=FitH, 
    linkcolor=linkcolor,
    urlcolor=urlcolor,
    colorlinks=true
}

\usepackage{import}

\title{\textbf{Элементы теории чисел}}
\author{Лектор: проф. Королев Максим Александрович}

\begin{document}
    
    \fontsize{14pt}{20pt}\selectfont
    \maketitle
    \vspace{0.3cm}
    \begin{center}
        \includegraphics[width=0.75\linewidth]{Pictures/mehmat.png}
    \end{center}
    \vspace{1.5cm}
    \begin{center}
        Конспект: Кирилл Яковлев, 108 группа\\
        Telegram: \href{https://t.me/fourkenz}{@fourkenz}\\
        GitHub: \href{https://github.com/yakovlevki}{yakovlevki}\\
    \end{center}
        
    \newpage
    \tableofcontents
    \fontsize{14pt}{20pt}\selectfont
    \newpage

    \import{Sections/}{section01.tex}
    \import{Sections/}{section02.tex}
    \import{Sections/}{section03.tex}
    \import{Sections/}{section04.tex}
    \import{Sections/}{section05.tex}
    \import{Sections/}{section06.tex}
    \import{Sections/}{section07.tex}
    \import{Sections/}{section08.tex}
    \import{Sections/}{section09.tex}
    \import{Sections/}{section10.tex}
    \import{Sections/}{section11.tex}

    \end{document}